\paragraph{Mechanism.}
On DKT, three independent pieces of evidence point to label copying: the gap between first and later KCs (Fig.~\ref{fig:kcpos}), the shuffle-by-question control (0.625 against 0.500--0.501) and the intervention test (1.00 against 0). On simpleKT only the gap is available, so the evidence is indirect.

\paragraph{Validation in pyKT.}
The KC-level validation AUC of pyKT exceeds the test AUC by about 0.05 for every model that leaks: $+0.052$ for DKT and for the pyKT DKT, $+0.046$ for AKT and $+0.049$ for simpleKT. For the models that satisfy (ii) and for QIKT the difference is between $-0.003$ and $-0.002$. Selecting a model with the pyKT validation would therefore favour a leaking one. This agrees with the inflated validation reported by Badran and Preisach~\cite{DBLP:conf/csedu/BadranP25a}.

\paragraph{Size of the effect.}
The overall effect is small ($+0.005$ and $+0.004$ AUC) because only 12.5\% of the questions have several KCs.


\paragraph{Limitations.}
We use one dataset; simpleKT has one seed; the shuffle, intervention and leak-free-selection controls ran on DKT only; no hyperparameter was tuned, and our pyKT baselines are 0.0016--0.0045 AUC below published values. The segment convention of simpleKT-PQC is the likely cause of its loss of 0.17 AUC on immediate repeats and of most of its loss on single-KC questions (Sect.~\ref{subsec:res-skt}); we have not isolated its effect, and a small residual loss on single-KC questions that are not repeats ($-$0.0008) cannot be excluded. DKT-PQC is 0.0067 below DKT on single-KC first graded questions, which we have not explained. We do not rank the three implementations of the rule.